% ===========================================================================
%  Clean formula subsections for the report:
%     (a) Wootters concurrence        (b) Werner state       (c) <S^2>
%  Drop-in replacements / additions. Nothing else in the report is affected.
% ===========================================================================

% ---------------------------------------------------------------------------
\subsection{Wootters Concurrence (bipartite, $n=2$)}
% ---------------------------------------------------------------------------
For a two-qubit mixed state $\rho$ (a $4\times4$ density matrix) the
entanglement is known in closed form~\cite{Wootters1998}. The construction
compares $\rho$ with its own \emph{spin-flipped} copy. For a pure state the
concurrence is simply the overlap with the flipped state,
$C=\lvert\langle\psi|\tilde\psi\rangle\rvert$, which vanishes for every product
state and equals unity for every maximally entangled state --- singlet,
triplet and Bell states alike. The flip therefore separates entangled from
separable states, as an entanglement measure must, without distinguishing
between the different maximally entangled states.

The flip is generated by the Pauli matrix
%
\begin{equation}
    \sigma_y = \begin{pmatrix} 0 & -i \\ i & 0 \end{pmatrix},
\end{equation}
%
which for two qubits becomes
%
\begin{equation}
    \mathcal{F} = \sigma_y \otimes \sigma_y =
    \begin{pmatrix}
        0 & 0 & 0 & -1 \\
        0 & 0 & 1 & 0 \\
        0 & 1 & 0 & 0 \\
        -1 & 0 & 0 & 0
    \end{pmatrix}.
\end{equation}
%
The spin-flipped state is obtained by the antiunitary transformation
%
\begin{equation}
    \tilde{\rho} = \mathcal{F}\, \rho^{*}\, \mathcal{F},
    \qquad
    R = \rho\,\tilde{\rho}
      = \rho\,\mathcal{F}\,\rho^{*}\,\mathcal{F},
\end{equation}
%
where $\rho^{*}$ is the complex conjugate of $\rho$. Note that $\mathcal{F}$
appears on \emph{both} sides of $\rho^{*}$: this conjugation (not a
projection) is reversible, preserves the trace and the spectrum, and returns a
valid density matrix, whereas $\mathcal{F}\rho^{*}$ alone would not. The
expression is quadratic in $\rho$, which enters once directly and once
conjugated.

Let $\mu_i$ be the eigenvalues of $R$ and $\lambda_i=\sqrt{\mu_i}$ their
square roots, sorted in decreasing order,
$\lambda_1\ge\lambda_2\ge\lambda_3\ge\lambda_4\ge0$. ($R$ is in general not
Hermitian, so a general eigenvalue routine is required.) The concurrence is
%
\begin{equation}
    \boxed{\;C(\rho) = \max\!\left(0,\;
           \lambda_1 - \lambda_2 - \lambda_3 - \lambda_4\right)\;}
\end{equation}
%
with $C=0$ for a separable state and $C=1$ for a maximally entangled (Bell)
state.

Wootters proved that this expression is the \emph{exact convex roof} for two
qubits. Consequently, for $n=2$ no numerical optimisation is required and
%
\begin{equation}
    C_{\mathrm{GME}}\big(\Gamma^{(2)}\big) \;=\; C\big(\Gamma^{(2)}\big).
\end{equation}
%
This provides an analytic benchmark for the numerical solver: the GCG convex
roof reproduces the Wootters value to machine precision (agreement of
$6\times10^{-16}$ for $L=6$ and $3\times10^{-15}$ for $L=8$ at $U=10$,
$t=0.05$).

% ---------------------------------------------------------------------------
\subsection{Werner State}
% ---------------------------------------------------------------------------
A Werner state has the form
%
\begin{equation}
    \rho_W = p\,\ket{\Psi^-}\bra{\Psi^-} + \frac{1-p}{4}\, I ,
\end{equation}
%
where $\ket{\Psi^-}=\tfrac{1}{\sqrt2}\bigl(\ket{\uparrow\downarrow}
-\ket{\downarrow\uparrow}\bigr)$ is the singlet and $p\in[0,1]$ is the mixing
parameter. Its eigenvalues are
%
\begin{equation}
    \lambda_{\max} = \frac{1+3p}{4},
    \qquad
    \lambda_2=\lambda_3=\lambda_4=\frac{1-p}{4},
\end{equation}
%
i.e.\ one non-degenerate eigenvalue and a three-fold degenerate one. This
degeneracy pattern is the signature of the Werner form and is reproduced
exactly by the computed two-site RDMs, which confirms that $\Gamma^{(2)}$ is
of Werner type. The parameter follows from the largest eigenvalue,
%
\begin{equation}
    p = \frac{4\lambda_{\max}-1}{3},
\end{equation}
%
and the concurrence reduces to
%
\begin{equation}
    C = \max\!\left(0,\ \frac{3p-1}{2}\right),
\end{equation}
%
in agreement with the general Wootters expression above.

A related quantity, obtained by mapping the concurrence through
%
\begin{equation}
    W(\rho) = \sqrt{\,1-\sqrt{1-C^{2}}\,},
\end{equation}
%
is also reported for reference. It should be noted that $W$ is a
\emph{different measure} from the concurrence: the two coincide only at the
extreme points $C=0$ and $C=1$. For the ground state at $U=10$, $t=0.05$ one
finds $C=0.4342$ but $W=0.3149$ ($L=6$), and $C=0.4127$ but $W=0.2985$
($L=8$). The convex-roof value quoted as $C_{\mathrm{GME}}$ throughout this
work is the concurrence $C$, consistent with the numerical solver.

% ---------------------------------------------------------------------------
\subsection{Total Spin $\langle S^{2}\rangle$ of a Cluster}
% ---------------------------------------------------------------------------
The spin character of an $n$-site cluster is quantified by the total
spin-squared operator
%
\begin{equation}
    \mathbf{S}^{2} = \Bigl(\sum_{i=1}^{n}\mathbf{S}_i\Bigr)^{2}
                   = \sum_{i=1}^{n}\sum_{j=1}^{n} \mathbf{S}_i\!\cdot\!\mathbf{S}_j ,
\end{equation}
%
where the spin-$\tfrac12$ operators act on the local spin basis
$\{\ket{\uparrow},\ket{\downarrow}\}$ and
%
\begin{equation}
    \mathbf{S}_i\!\cdot\!\mathbf{S}_j
      = S_i^{z}S_j^{z}
      + \tfrac{1}{2}\bigl(S_i^{+}S_j^{-} + S_i^{-}S_j^{+}\bigr).
\end{equation}
%
For a state of definite total spin $S$ one has
$\mathbf{S}^{2}\ket{\psi}=S(S+1)\ket{\psi}$, so that
%
\begin{equation}
    S=0 \Rightarrow \langle S^{2}\rangle = 0, \qquad
    S=\tfrac12 \Rightarrow \langle S^{2}\rangle = \tfrac34, \qquad
    S=1 \Rightarrow \langle S^{2}\rangle = 2 .
\end{equation}
%
For the reduced density matrix of a cluster the expectation value is
%
\begin{equation}
    \boxed{\;\langle S^{2}\rangle
        = \mathrm{Tr}\bigl[\Gamma^{(n)}\,\mathbf{S}^{2}\bigr]
        = \sum_{S} p_S\, S(S+1)\;}
\end{equation}
%
where $p_S$ is the weight of total-spin sector $S$. Although $\mathbf{S}^{2}$
commutes with $H$ and total spin is therefore a good quantum number of the
\emph{full} ring, it is not conserved for a \emph{subsystem}: the cluster
exchanges spin with the remaining $L-n$ sites and is generally in a mixture of
spin sectors. For this reason $\langle S^{2}\rangle$ is reported directly and
no effective $S$ is extracted, since inverting $S(S+1)$ would yield a
non-integer, non-physical value.

The sector decomposition makes the spin content explicit. For $L=6$,
$U=10$, $t=0.05$ at $T=0$ one obtains
%
\begin{equation}
    n=2:\; 71.7\%\ \text{singlet} + 28.3\%\ \text{triplet}
    \;\Rightarrow\; \langle S^{2}\rangle = 0.5658,
\end{equation}
\begin{equation}
    n=3:\; 98.4\%\ \text{doublet} + 1.6\%\ \text{quartet}
    \;\Rightarrow\; \langle S^{2}\rangle = 0.7975 .
\end{equation}
%
An odd-sized cluster cannot pair all of its spins internally and therefore
carries at least $S=\tfrac12$, i.e.\ $\langle S^{2}\rangle\ge\tfrac34$, whereas
an even-sized cluster can be fully paired and may reach
$\langle S^{2}\rangle=0$. This even/odd alternation is the origin of the
non-monotonic behaviour of $C_{\mathrm{GME}}$ with cluster size.
